\par\noindent\textbf{(1)} 线性化质量守恒为 \(\eta_t+\partial_x(hu)=0\)，水平动量为 \(u_t+g\eta_x=0\)。当水深 \(h\) 常数，消去 \(u\) 得 \(\eta_{tt}=gh\,\eta_{xx}\)，是传播速度 \(\sqrt{gh}\) 的双向波动方程。
\par\noindent\textbf{(2)} 设 \(\eta\propto e^{i(kx-\omega t)}\)，得到 \(\omega^2=ghk^2\)，故相速度与群速度都为 \(c=\sqrt{gh}\)。平均单位面积总波能密度约为 \(E=\rho g\langle\eta^2\rangle\)，能流为 \(F=Ec\)，需同时计入等量的动能和重力势能。
\par\noindent\textbf{(3)} 若 \(h(x)\) 在单波长内变化缓慢，波作用守恒近似为平均能流 \(F=\rho g\langle\eta^2\rangle\sqrt{gh}\) 沿传播近似不变。于是振幅 \(A\propto h^{-1/4}\)，称Green定律；海岸变浅时波速降低、振幅升高，并非简单按 \(h^{-1/2}\) 变化。
\par\noindent\textbf{(4)} 当浅水条件 \(kh\ll1\) 不满足，完整水波色散应为 \(\omega^2=gk\tanh(kh)\)。近岸破碎、摩擦、反射及非线性陡化也破坏无耗散缓变假设，因此不能把 \(h^{-1/4}\) 放大律延伸到 \(h\to0\) 的岸线。
